\section{Prime--M\"obius leg fusion and the joint-mask boundary}
\label{sec:tpc32-prime-mobius}

The factorization of the cutoff shell becomes useful only if the actual
opened-row coefficient is retained.  We now restore that provenance and
separate what is algebraically exact from what would require a new norm
estimate.

\subsection{The physical row coefficient}

For $i\in\{1,2\}$, a row is $\alpha=(\ell,d)$ with
\begin{equation}
 \ell\asymp L\ \text{prime},
 \qquad
 d\asymp D\ \text{squarefree},
 \qquad
 (d,H)=1,
 \qquad
 m_\alpha=\ell d\asymp Q,
 \label{eq:tpc32-row-provenance-fusion}
\end{equation}
and has coefficient
\begin{equation}
 \gamma_\alpha^{(i)}
 =\mu(d)(\log\ell)\omega_D^{(i)}(d)
   \psi_L^{(i)}(\ell/L)\zeta_\alpha^{(i)}.
 \label{eq:tpc32-gamma}
\end{equation}
Here the $\omega_D^{(i)}$ and $\psi_L^{(i)}$ range over fixed bounded
smooth families and $\zeta_\alpha^{(i)}$ contains fixed-period characters
or phases.  The inherited estimates are
\begin{equation}
 \sum_\alpha|\gamma_\alpha^{(i)}|\ll Q,
 \qquad
 \sum_\alpha|\gamma_\alpha^{(i)}|^2\ll Q\log(2L).
 \label{eq:tpc32-row-masses}
\end{equation}
These bounds use Chebyshev's estimate for the prime source and are stronger
than replacing every logarithm by $X^\eps$ \citep{WangTPC25}.

Recall the squarefree divisor coefficient
\begin{equation}
 b_R(w)=-\mu(w)w_R(w),
 \label{eq:tpc32-bR-squarefree}
\end{equation}
where
\begin{equation}
 w_R(w)=\log w+\one_{w\le R}\frac{w}{\varphi(w)}
 \sum_{\substack{a\le R/w\\(a,w)=1}}
 \frac{\mu(a)^2}{\varphi(a)}
 \label{eq:tpc32-wR-fusion}
\end{equation}
on squarefree $w$; both sides of \eqref{eq:tpc32-bR-squarefree} vanish off
the inherited squarefree support.  For $w>T\ge R$ this reduces to
\begin{equation}
 b_R(w)=a(w)=-\mu(w)\log w.
 \label{eq:tpc32-ultra-a}
\end{equation}

\begin{lemma}[Primitive one-leg sign fusion]
\label{lem:tpc32-leg-fusion}
Suppose that $w\mid m_\alpha j+h$ on the primitive physical support.
Then $(w,d)=1$, and
\begin{equation}
 \boxed{
 \gamma_\alpha^{(i)}b_R(w)
 =-\mu(dw)(\log\ell)\omega_D^{(i)}(d)
   \psi_L^{(i)}(\ell/L)\zeta_\alpha^{(i)}w_R(w).}
 \label{eq:tpc32-leg-fusion}
\end{equation}
In the ultra range $w>T$,
\begin{equation}
 \mu(d)a(w)=-\mu(dw)\log w.
 \label{eq:tpc32-ultra-leg-fusion}
\end{equation}
\end{lemma}

\begin{proof}
If a prime $p$ divided both $w$ and $d$, then it would divide
$m_\alpha=\ell d$ and $m_\alpha j+h$, hence $p\mid h$.  This contradicts
$(d,H)=1$, because $H=\operatorname{rad}|h|$.  Thus $(w,d)=1$.
The divisor coefficient is squarefree-supported, and $d$ is squarefree, so
multiplicativity gives $\mu(d)\mu(w)=\mu(dw)$.  Substitute
\eqref{eq:tpc32-bR-squarefree} into \eqref{eq:tpc32-gamma};
\eqref{eq:tpc32-ultra-leg-fusion} follows from
\eqref{eq:tpc32-ultra-a}.
\end{proof}

One may consequently write a cutoff row signal schematically as
\begin{align}
 \cT_W^{(i)}(j)
 ={}&-\sum_{\ell,d}\sum_{\substack{w\le W\\w\mid\ell dj+h}}
 \mu(dw)(\log\ell)w_R(w)
 \omega_D^{(i)}(d)\psi_L^{(i)}(\ell/L)
 \zeta_{\ell,d}^{(i)}\Phi_{\ell,d}^{(i)}(j)
 \notag\\
 &\quad+\text{the calibrated drift},
 \label{eq:tpc32-one-row-signal}
\end{align}
where $\Phi$ retains the actual fixed-residue and smooth physical factors.
Changing variables to $n=dw$ makes the outer sign a single $\mu(n)$, but
the inner coefficient still contains the factorization $n=dw$, the prime
source $\ell$, and the congruence $w\mid\ell dj+h$.  Thus
\eqref{eq:tpc32-one-row-signal} is a prime--M\"obius normal form, not a
claim that the resulting coefficient is well-factorable in the technical
sieve-theoretic sense.

\subsection{What the smooth and residue factors permit}

Fixed-period masks have a finite Fourier expansion of bounded total mass.
The smooth factors have an absolutely convergent Mellin projective
representation; after the actual row coefficients have been included, its
weighted projective $\ell^1$ mass is $O(Q^2)$ and its energy mass is
$O(Q\log(2L))$ \citep{WangTPC25}.  Consequently these specific factors may
be separated and reassembled without introducing an $X^\eps$ projective
loss.

The generic row-pair mask is different.  Write the full joint multiplier as
\begin{equation}
 \mathfrak A_{\alpha,\gamma}(j)
 =\mathfrak m(\alpha,\gamma)
   \Xi_{\alpha,\gamma}(j)
   \cW_{\alpha,\gamma}(j/J).
 \label{eq:tpc32-joint-multiplier-boundary}
\end{equation}
The inherited abstract interface assumes only
\begin{equation}
 |\mathfrak m(\alpha,\gamma)|\le1,
 \qquad
 \mathfrak m(\alpha,\gamma)=0
 \quad\text{when }m_\alpha=m_\gamma.
 \label{eq:tpc32-mask-bound}
\end{equation}
Entrywise multiplication by such a matrix need not preserve positivity,
rank, or spectral norm.  Accordingly, the exact formula in
\cref{sec:tpc32-canonical-factorization} is a factorable row core with
\eqref{eq:tpc32-joint-multiplier} retained; it is not an unqualified product
of two row sums.

\begin{definition}[Controlled projective multiplier]
\label{def:tpc32-projective-multiplier}
We say that $\mathfrak A$ has projective cost $\mathfrak P$ on a fixed
physical packet if it admits a representation
\begin{equation}
 \mathfrak A_{\alpha,\gamma}(j)
 =\int_\Omega p_{\omega,j}(\alpha)q_{\omega,j}(\gamma)
 \,d\nu(\omega)
 \label{eq:tpc32-projective-decomposition}
\end{equation}
whose required row norms and orbit seminorms integrate to at most
$\mathfrak P$.
\end{definition}

Under \eqref{eq:tpc32-projective-decomposition}, each bilinear core may be
written as an integral of products of row sums and any estimate uniform in
the displayed seminorms may be reassembled with cost $\mathfrak P$.
Finite-dimensional singular-value decomposition alone does not give the
needed theorem: its projective norm may be as large as the obstacle scale.
No controlled decomposition for the generic pair mask is asserted here.

\begin{remark}[Two different factorization statements]
The cutoff coefficient has algebraic rank at most two in the divisor
variables.  The pair mask is a Schur multiplier in the row variables.
These are independent statements.  Low divisor rank does not remove an
arbitrary joint row multiplier, and bounded entries do not imply bounded
Schur or projective norm.
\end{remark}
